\section{Modelos del sistema}
\label{sec:modelos}

Siguiendo a Coulouris et al.~\cite{coulouris}, el sistema se describe con cuatro
modelos. Para cada uno se presentan los conceptos y cómo se aplican al proyecto.

% =====================================================================
\subsection{Modelo arquitectónico}

El modelo arquitectónico define qué entidades forman el sistema, cómo se comunican,
qué papel cumple cada una y dónde se ubican.

\subsubsection{Estilo arquitectónico}

\begin{table}[H]
\centering
\caption{Estilos arquitectónicos aplicados al proyecto.}
\begin{tabularx}{\textwidth}{>{\bfseries}p{3.6cm}Y}
\toprule
Concepto & Cómo se aplica \\
\midrule
Arquitectura por capas &
Tres capas en orden estricto: \emph{Consumidores} $\rightarrow$ \emph{Servicios de
Negocio} $\rightarrow$ \emph{Servicios de Almacenamiento}. Cada capa solo llama a la
capa inmediatamente inferior, por lo que cada una puede desplegarse, reemplazarse y
escalarse por separado. \\
Cliente--servidor &
Los clientes son clientes del \codigo{GestorTransacciones}; a su vez, los procesos
de negocio son clientes del Servicio de Persistencia. \\
Intermediario (broker) &
El \codigo{GestorTransacciones} es el único punto de entrada: los clientes solo
conocen su dirección y él reparte las peticiones entre los procesos de negocio. Los
clientes quedan desacoplados de cuántos procesadores hay y dónde están. \\
Servicios sin estado &
Gestor, procesadores y servicio de cancelación no guardan datos de negocio en
memoria; todo el estado vive en la capa de almacenamiento. Por eso pueden
reiniciarse o replicarse sin perder información. \\
Publicación--suscripción &
Las cancelaciones se anuncian con ZeroMQ \codigo{PUB}/\codigo{SUB}: el gestor
publica sin conocer a los suscriptores (desacoplamiento en el espacio) y la
solicitud guardada permite aplicarla después (desacoplamiento en el tiempo). \\
Fuente única de verdad &
El Servicio de Persistencia es el único que decide los cupos, las reservas y las
transacciones, lo que simplifica la consistencia. \\
\bottomrule
\end{tabularx}
\end{table}

\subsubsection{Entidades y roles}

Las entidades corresponden a las cajas de la Figura~\ref{fig:arquitectura-basica}.
Cada una es un proceso independiente:

\begin{table}[H]
\centering
\caption{Procesos del sistema.}
\begin{tabularx}{\textwidth}{p{4.1cm}p{2.6cm}Y}
\toprule
Proceso & Capa & Responsabilidad \\
\midrule
\codigo{cliente} & Consumidores & Interfaz del estudiante: consultar eventos del mes,
reservar, modificar y cancelar; recibe notificaciones de cancelación. \\
\codigo{gestor-transacciones} & Servicios de Negocio & Recibe todas las solicitudes,
envía las síncronas a un procesador RMC libre y publica las cancelaciones. \\
\codigo{procesador-rmc} & Servicios de Negocio & Procesamiento de
\textbf{R}eservas, \textbf{M}odificaciones y \textbf{C}onsultas: valida y ejecuta
las operaciones síncronas contra persistencia. \\
\codigo{cancelacion-reservas} & Servicios de Negocio & Aplica de forma asíncrona las
cancelaciones aceptadas y notifica el resultado al estudiante. \\
\codigo{persistencia} & Servicios de Almacenamiento & Guarda estudiantes, lugares,
eventos, fechas, reservas y transacciones; aplica cada cambio de forma atómica. \\
\bottomrule
\end{tabularx}
\end{table}

\subsubsection{Paradigmas de comunicación}

Toda la comunicación entre procesos usa paso de mensajes con ZeroMQ~\cite{zguide}:

\begin{itemize}
    \item \textbf{Petición--respuesta} (\codigo{REQ}/\codigo{ROUTER}/\codigo{DEALER}/\codigo{REP})
          para consultas, reservas y modificaciones, y para la aceptación de la
          cancelación. Quien llama se bloquea hasta recibir la respuesta.
    \item \textbf{Publicación--suscripción} (\codigo{PUB}/\codigo{SUB}) para anunciar
          una cancelación aceptada al servicio de cancelación y para notificar el
          resultado al estudiante. El publicador no espera.
\end{itemize}

\subsubsection{Ubicación física}

En la arquitectura base cada proceso corre en una máquina del laboratorio
(sección~\ref{sec:despliegue}): una máquina para los clientes, una para el
gestor, una para los procesos de negocio y una para persistencia. Las direcciones y
puertos se leen de variables de entorno, de modo que el mismo programa corre en una
sola máquina (desarrollo) o repartido en varias (pruebas). En la segunda entrega el
sistema se desplegará en Kubernetes sobre las mismas 4 máquinas
(sección~\ref{sec:escalabilidad}).

% =====================================================================
\subsection{Modelo de interacción}

El modelo de interacción trata el desempeño de los canales, la ausencia de un reloj
global y el orden de los eventos.

\subsubsection{Sistema asíncrono y cotas prácticas}

En un \emph{sistema distribuido síncrono} existen cotas conocidas para el tiempo de
ejecución de los procesos, el retardo de los mensajes y la deriva de los relojes;
en uno \emph{asíncrono} no existen. La red del laboratorio no garantiza ninguna
cota, así que el sistema corre sobre un modelo \textbf{asíncrono}. Que las
operaciones sean ``síncronas'' se refiere al \emph{estilo de interacción} (el
cliente espera la respuesta), no a tiempos garantizados.

Para poder actuar, el sistema asume \textbf{cotas prácticas} mediante timeouts:
\begin{itemize}
    \item cliente $\rightarrow$ gestor: 1\,s por intento, hasta 3 intentos;
    \item procesos de negocio $\rightarrow$ persistencia: 300\,ms por intento, hasta 3 intentos;
    \item latidos (\emph{heartbeats}) entre gestor y procesadores cada 1\,s; tres
          perdidos significan que el procesador está caído.
\end{itemize}
Cuando se vence un timeout el resultado se considera \textbf{desconocido}, no
fallido: la operación pudo haberse aplicado. Por eso los reintentos deben ser
idempotentes (sección~\ref{sec:semantica}).

\subsubsection{Desempeño de los canales}

\begin{itemize}
    \item \textbf{Latencia:} ZeroMQ mantiene conexiones TCP persistentes entre
          procesos, así que no se paga un establecimiento de conexión por petición.
    \item \textbf{Ancho de banda:} los mensajes son pequeños (texto
          \codigo{clave=valor}); la consulta de un mes devuelve pocas decenas de filas.
    \item \textbf{Jitter:} se observa comparando el promedio con la desviación
          estándar y el percentil 95 en las pruebas de la Tabla~1.
\end{itemize}

\subsubsection{Relojes y orden de los eventos}

Cada máquina tiene su propio reloj y estos derivan. Por eso:
\begin{itemize}
    \item Las \textbf{duraciones} (tiempos de respuesta) se miden siempre con el reloj
          monótono de una sola máquina: inicio y fin en el mismo proceso cliente.
    \item El \textbf{orden de las operaciones} sobre los mismos datos lo decide el
          Servicio de Persistencia, que atiende las peticiones una a una. Si dos
          estudiantes compiten por el último cupo, gana el que persistencia procesa
          primero, no el que tenga la marca de tiempo menor.
    \item Las marcas de tiempo guardadas (instante de una transacción) las pone
          siempre persistencia, así que son consistentes entre sí.
    \item Cada operación lleva un \codigo{idSolicitud} único que permite seguirla a
          través de los procesos en los registros.
\end{itemize}

\subsubsection{Semántica de invocación}
\label{sec:semantica}

\begin{table}[H]
\centering
\caption{Semántica de invocación por operación.}
\begin{tabularx}{\textwidth}{p{4.2cm}p{2.6cm}Y}
\toprule
Operación & Estilo & Semántica \\
\midrule
Consultar eventos del mes & Síncrona & \emph{Al menos una vez}: los reintentos son
inofensivos porque es de solo lectura. \\
Reservar, modificar & Síncrona & \emph{Efectivamente exactamente una vez}: el
cliente reintenta con el mismo \codigo{idSolicitud} y persistencia devuelve la
respuesta guardada sin aplicar el cambio otra vez. \\
Cancelar (aceptación) & Síncrona, corta & Exactamente una solicitud de cancelación
guardada por \codigo{idSolicitud}. \\
Cancelar (aplicación) & Asíncrona & Entrega \emph{al menos una vez} (aviso
PUB/SUB + barrido periódico) y aplicación idempotente: la cancelación se aplica una
sola vez. \\
\bottomrule
\end{tabularx}
\end{table}

% =====================================================================
\subsection{Modelo de fallos}

El modelo de fallos define qué puede fallar, cómo falla y cómo el sistema lo detecta
y lo enmascara.

\subsubsection{Clases de fallos considerados}

\begin{table}[H]
\centering
\caption{Clases de fallos.}
\begin{tabularx}{\textwidth}{p{3.3cm}Y Y}
\toprule
Clase & Definición & Ejemplo en el proyecto \\
\midrule
Caída (\emph{fail-stop}) & El proceso se detiene y no hace nada más. &
Se apaga la máquina del procesador RMC. \\
Omisión de proceso & No envía o no recibe un mensaje. & El servicio de cancelación
se cae justo después de recibir un aviso. \\
Omisión de canal & El mensaje se pierde en la red. & Se pierde la respuesta de
\codigo{reservar}; el cliente agota el timeout. \\
Temporización & La respuesta llega, pero tarde. & Persistencia tarda más de 1\,s
bajo carga. \\
Arbitrario (bizantino) & Comportamiento incorrecto o malicioso. & Mensaje alterado
o cliente que envía datos inválidos. Queda fuera del modelo de corrección y lo
cubre en parte el modelo de seguridad. \\
\bottomrule
\end{tabularx}
\end{table}

\textbf{Supuestos.} (1) Los procesos fallan por caída. (2) Los canales pueden
perder, retrasar o duplicar mensajes, pero TCP detecta la corrupción. (3) Lo que
persistencia confirma como guardado sobrevive a una caída; los mensajes de ZeroMQ,
incluidos los de PUB/SUB, no son durables. (4) Desde afuera no se distingue una
caída de un proceso lento: un timeout solo significa ``desconocido''.

\subsubsection{Fallos por componente y cómo se enmascaran}

\begin{longtable}{p{3.1cm}p{3.6cm}p{8.6cm}}
\caption{Fallos por componente.}\\
\toprule
Componente & Fallo & Detección y manejo \\
\midrule
\endfirsthead
\toprule
Componente & Fallo & Detección y manejo \\
\midrule
\endhead
Red cliente--gestor & Se pierde la petición o la respuesta & Timeout del cliente;
reintento con el mismo \codigo{idSolicitud} (patrón \emph{Lazy Pirate}); no hay
reservas duplicadas. \\
Procesador RMC & Se cae con una petición en curso & El gestor deja de recibir sus
latidos y lo descarta; el cliente reintenta y otro procesador atiende. \\
Gestor de Transacciones & Se cae & Punto único de fallo en la arquitectura base; los
clientes reintentan hasta que se reinicie. Al volver, los procesadores se
registran de nuevo con su siguiente latido. En la segunda entrega se replica
(\emph{Binary Star}). \\
Servicio de cancelación & Se cae o pierde un aviso PUB/SUB & La solicitud ya quedó
guardada como \codigo{PENDIENTE}; un barrido periódico la encuentra y la aplica. \\
Persistencia & Se cae a mitad de una operación & Cada cambio es atómico (todo o
nada) y se escribe en disco antes de responder; al reiniciar se reconstruye el
estado. Los servicios responden \codigo{PERSISTENCIA\_NO\_DISPONIBLE}. \\
Persistencia & Lenta & Los indicadores de latencia lo muestran (p95 > 1\,s). \\
Relojes & Deriva entre máquinas & Las duraciones se miden en una sola máquina y el
orden lo decide persistencia. \\
\bottomrule
\end{longtable}

\subsubsection{Puntos únicos de fallo de la arquitectura base}

En la arquitectura base hay una sola instancia de cada componente, así que el
gestor, el servicio de cancelación y persistencia son puntos únicos de fallo. La
sección~\ref{sec:escalabilidad} presenta los componentes que los eliminan.

% =====================================================================
\subsection{Modelo de seguridad}

El modelo de seguridad identifica qué se protege, de quién y con qué defensas. Se
asume un \emph{enemigo} capaz de enviar cualquier mensaje a cualquier proceso y de
leer, copiar, alterar o repetir los mensajes que viajan por la red.

\subsubsection{Objetos a proteger}

\begin{itemize}
    \item \textbf{Datos personales del estudiante}: nombre, tipo y número de
          documento, correo, teléfono, carrera, departamento y semestre. En Colombia
          están protegidos por la Ley 1581 de 2012~\cite{ley1581}.
    \item \textbf{Reservas y cupos}: recurso limitado que podría acapararse.
    \item \textbf{Transacciones}: registro de auditoría de reservas, modificaciones y
          cancelaciones, con sus montos.
    \item \textbf{Eventos y lugares}: datos falsos engañarían a todos los estudiantes.
\end{itemize}

\subsubsection{Amenazas y defensas}

\begin{longtable}{p{3.0cm}p{5.0cm}p{7.3cm}}
\caption{Amenazas y defensas.}\\
\toprule
Amenaza & Ejemplo & Defensa \\
\midrule
\endfirsthead
\toprule
Amenaza & Ejemplo & Defensa \\
\midrule
\endhead
Suplantación & Cancelar la reserva de otro estudiante. & Autenticación del
estudiante (token de sesión) y verificación de propiedad: el
\codigo{idEstudiante} sale del token, nunca del cuerpo del mensaje
(\codigo{RESERVA\_AJENA}). Procesos internos autenticados con llaves CurveZMQ. \\
Escucha & Leer números de documento en la red. & Cifrado de todos los sockets con
CurveZMQ. \\
Alteración & Cambiar \codigo{cupos=1} por \codigo{cupos=50}. & CurveZMQ (cifrado
autenticado) y validación de rangos en el procesador RMC. \\
Repetición & Reenviar una reserva capturada. & \codigo{idSolicitud} único: la
repetición devuelve la respuesta guardada sin un segundo efecto. \\
Denegación de servicio & Inundar \codigo{reservar}. & Límite de solicitudes por
estudiante, colas acotadas de ZeroMQ (\emph{high-water marks}). \\
Elevación de privilegios & Un estudiante crea o borra eventos. & Autorización por
rol en la capa de negocio. \\
Fuga de datos & Documentos en los registros. & Los registros solo llevan
\codigo{idSolicitud} e \codigo{idEstudiante}; el documento se muestra enmascarado. \\
Repudio & ``Yo nunca cancelé''. & Toda operación queda como \codigo{Transaccion}
inmutable con estudiante, solicitud e instante. \\
Aislamiento & Conectarse directo a persistencia. & Solo el gestor es accesible desde
fuera; persistencia solo acepta conexiones de los procesos de negocio
(firewall o \emph{NetworkPolicy} en Kubernetes). \\
\bottomrule
\end{longtable}

\subsubsection{Seguridad frente a desempeño}

El cifrado agrega costo a cada mensaje. Como ZeroMQ mantiene conexiones
persistentes, el intercambio de llaves de CurveZMQ ocurre una sola vez por conexión
y no en cada petición. El costo se medirá comparando corridas con y sin cifrado.
En la primera entrega la autenticación y el cifrado no están implementados; se
incorporarán con los componentes de la sección~\ref{sec:escalabilidad}.
